% ============================================================
% Section: Experimental Setup
% (structure follows the reference paper: 4.1 environment & training,
%  4.2 evaluation protocol & metrics with A/B/C sub-parts)
% Preamble needs:  \usepackage{booktabs,xcolor,soul,amsmath,amssymb}
% \tbd{...} = information still to be filled in (yellow highlight)
% ============================================================


\section{Experimental Setup}

This section presents the experimental setup used to evaluate the proposed
framework for automated estimation of root canal filling length. Specifically,
the hardware and software configurations, training settings,
and hyperparameter selections are introduced to ensure experimental
reproducibility. The evaluation protocol and the performance metrics used for
tooth segmentation, length estimation, and clinical safety are then defined.

\subsection{Implementation Environment and Training Configuration}

All experiments were conducted on a local workstation running Microsoft
Windows~11 (64-bit). The hardware comprised an AMD Ryzen~7 9700X processor
(8 cores, 3.80~GHz), 32~GB of RAM, and an NVIDIA GeForce RTX~5070~Ti GPU with
16~GB of VRAM to support deep learning model training and inference. The
software stack was implemented in Python~3.14.6, with PyTorch~2.13.0 and
CUDA~13.0 (cuDNN~9.20) as the deep learning framework and GPU acceleration
backend. Ultralytics~8.4.104 was used for training and inference of the
detection and segmentation models, and OpenCV~5.0.0 was used for image
preprocessing, mask post-processing, and geometric measurement.
Pixel-to-millimetre conversion and statistical analyses were performed in
MATLAB~R2026a with the Statistics and Machine Learning Toolbox and the Deep
Learning Toolbox. The detailed hardware and software configurations are
summarised in Table~\ref{tab:hw}.

\begin{table}[t]
\centering
\caption{Hardware and software configurations.}
\label{tab:hw}
\small
\begin{tabular}{@{}p{0.36\columnwidth}p{0.58\columnwidth}@{}}
\toprule
Component & Specification / Version \\
\midrule
Operating system & Microsoft Windows 11 (64-bit, build 26200) \\
GPU & NVIDIA GeForce RTX 5070 Ti (16 GB VRAM) \\
CPU / RAM & AMD Ryzen 7 9700X (8 cores, 3.80 GHz) / 32 GB (5600 MT/s) \\
Programming language & Python 3.14.6 \\
Deep learning framework & PyTorch 2.13.0 / CUDA 13.0 (cuDNN 9.20) \\
Detection / segmentation framework & Ultralytics 8.4.104 \\
Image analysis tool & OpenCV 5.0.0 \\
Numerical computing & MATLAB R2026a (Statistics and ML Toolbox, Deep Learning Toolbox) \\
\bottomrule
\end{tabular}
\end{table}


The hyperparameter settings are summarised in Table~\ref{tab:hyper}. The YOLO11n detector used for single-tooth extraction was initialised from \texttt{yolo11n.pt} and trained on the 179 full radiographs; its settings are also listed in Table~\ref{tab:hyper}. The
segmentation model was initialised from the pretrained \texttt{yolo11n-seg.pt}
weights and optimised with AdamW at an initial learning rate of
$1\times10^{-3}$. The model was trained for 150 epochs with a batch size of 4,
with early stopping disabled (patience $=0$) and automatic mixed precision
turned off to avoid numerical instability; the checkpoint with the best
validation fitness was retained. Data augmentation comprised random rotation
($\pm10^\circ$), translation (0.05), scaling (0.1), and mosaic (0.3), whereas
horizontal flipping was disabled to preserve the physiological orientation of
the teeth. The optimisation objective followed the standard segmentation losses
of the YOLO framework, including bounding-box regression, classification, and
mask losses. A fixed random seed (\tbd{}) and deterministic mode were used for reproducibility.

For pixel-to-millimetre conversion, Gaussian process regression (GPR; MATLAB
\texttt{fitrgp}) took the long-side-normalised pixel length and patient age as
inputs. A squared-exponential kernel with a linear basis (mean) function was
adopted, and the kernel hyperparameters (length scale, signal standard
deviation, and noise standard deviation) were estimated within each training
fold by maximising the marginal log-likelihood. Input features were z-score
standardised with the mean and standard deviation of the training fold before
being applied to the test fold. Performance was assessed on out-of-fold
predictions from five-fold cross-validation (the 67 teeth were randomly permuted with a fixed seed of 42 and assigned to the five folds in turn, giving 13--14 teeth per fold without stratification), and all comparison settings---alternative input
features and ordinary least-squares linear regression with identical
inputs---used the same folds.
Because GPR involves no random initialisation, its results are deterministic for
a fixed fold assignment and were therefore obtained from a single run.

Comparative experiments followed a one-factor-at-a-time design, with the input
features fixed to the long-side-normalised pixel length and age. For the
preprocessing comparison, GPR was fixed and top-hat enhancement was compared with
CLAHE (clip limit \tbd{}, $8\times8$ tiles) and with the unprocessed images; a
segmentation model was trained for each group under identical training settings,
using the same polygon annotations and the same folds. For the regression
comparison, CLAHE preprocessing was fixed and GPR was compared with a
feed-forward artificial neural network (ANN; MATLAB \texttt{fitnet}, one hidden
layer of five neurons, Bayesian regularisation \texttt{trainbr}, at most 200
epochs); because of its random weight initialisation, ANN results are averaged
over 10 repetitions. Ordinary least-squares linear regression with identical
inputs served as an additional baseline.

\begin{table*}[t]
\centering
\caption{Summary of hyperparameter settings for model training.}
\label{tab:hyper}
\begin{tabular}{@{}lp{0.30\linewidth}p{0.44\linewidth}@{}}
\toprule
Model & Parameter & Value / Setting \\
\midrule
YOLO11n detector & Pretrained weights  & \texttt{yolo11n.pt} \\
            & Data split                & 179 radiographs; \tbd{} train / \tbd{} val \\
            & Input size                & \tbd{} \\
            & Epochs / batch size       & \tbd{} / \tbd{} \\
            & Optimiser / initial LR    & \tbd{} / \tbd{} \\
\midrule
YOLO11n-seg & Pretrained weights        & \texttt{yolo11n-seg.pt} \\
            & Input size                & $640\times640$ (letterbox) \\
            & Optimiser                 & AdamW \\
            & Initial learning rate     & $1\times10^{-3}$ \\
            & Batch size                & 4 \\
            & Epochs                    & 150 \\
            & Early stopping            & Disabled (patience $=0$) \\
            & Mixed precision (AMP)     & Off \\
            & Augmentation              & Rotation $\pm10^\circ$, translation 0.05, scale 0.1, mosaic 0.3, no flipping \\
            & Random seed               & \tbd{} \\
\midrule
GPR         & Inputs                    & Normalised pixel length + age (z-scored per training fold) \\
            & Kernel / basis            & Squared exponential / linear \\
            & Hyperparameter estimation & Maximum marginal log-likelihood, per fold \\
            & Validation                & 5-fold CV, out-of-fold ($n=67$) \\
            & Safety offset / tolerance & $-0.5$~mm / $\pm1.0$~mm \\
\midrule
ANN (comp.) & Architecture           & Feed-forward, 1 hidden layer, 5 neurons (\texttt{fitnet}) \\
            & Training algorithm        & Bayesian regularisation (\texttt{trainbr}), $\le$200 epochs \\
            & Repetitions               & 10 (random initialisation), mean reported \\
\bottomrule
\end{tabular}
\end{table*}

\subsection{Evaluation Protocol and Performance Metrics}

To quantitatively evaluate the proposed framework, performance was assessed at
three levels: tooth segmentation, root canal filling length estimation, and
clinical safety. Segmentation was evaluated on the validation and independent
test sets, whereas all length-related metrics were computed on out-of-fold
predictions after subtracting the 0.5~mm safety offset.

\paragraph{A. Segmentation Metrics}
Segmentation performance was reported as box and mask mAP@50 and mAP@50:95.
The mAP@50 metric evaluates the mean average precision at an IoU threshold of
0.50, while mAP@50:95 averages it over IoU thresholds from 0.50 to 0.95 in
steps of 0.05. On the independent test set, per-image mask mAP@50:95 was
computed at the original image resolution from all predicted masks with a
confidence of at least 0.001.

\paragraph{B. Length Estimation Metrics}
Let $y_i$ denote the reference filling length of the $i$-th tooth, $\hat{y}_i$
the offset-corrected predicted length, and $n$ the number of test teeth.
Estimation accuracy was evaluated by the mean absolute error (MAE), root mean
square error (RMSE), Pearson correlation coefficient ($R$), and coefficient of
determination ($R^2$):
\begin{align}
\mathrm{MAE}  &= \frac{1}{n}\sum_{i=1}^{n}\lvert \hat{y}_i - y_i \rvert, \\
\mathrm{RMSE} &= \sqrt{\frac{1}{n}\sum_{i=1}^{n}(\hat{y}_i - y_i)^2}, \\
R &= \frac{\sum_{i=1}^{n}(\hat{y}_i-\bar{\hat{y}})(y_i-\bar{y})}
         {\sqrt{\sum_{i=1}^{n}(\hat{y}_i-\bar{\hat{y}})^2}\sqrt{\sum_{i=1}^{n}(y_i-\bar{y})^2}}, \\
R^2 &= (R)^2 .
\end{align}
Because the safety offset only shifts the predictions, $R$ and $R^2$ are
unaffected by it, whereas MAE and RMSE change.

\paragraph{C. Clinical Safety Metrics}
Clinical safety was assessed by the ideal rate, defined as the proportion of
teeth whose corrected error falls within $\pm1.0$~mm, and by the mean bias,
where a positive value indicates systematic overestimation:
\begin{align}
\text{Ideal rate} &= \frac{1}{n}\sum_{i=1}^{n}\mathbb{1}\left(\lvert \hat{y}_i - y_i \rvert \le 1.0\right)\times 100\%, \\
\text{Bias} &= \frac{1}{n}\sum_{i=1}^{n}(\hat{y}_i - y_i).
\end{align}
The association between estimation error and age was examined with the
Spearman rank correlation, and differences between sexes were tested with the
Wilcoxon rank-sum test.
